\subsection{Поперечные моды ОВ}\label{subsec:modes}

Рассмотрим волновое уравнение~\eqref{eq:wave} в отсутствие нелинейной поляризации ($\mathbf{P_{\mathrm{NL}}}=0$). Данное приближение допустимо, поскольку в кварцевых ОВ нелинейная восприимчивость третьего порядка $\chi^{(3)}$ много меньше линейной восприимчивости $\chi^{(1)}$, и при не слишком высоких напряжённостях электрического поля нелинейным вкладом в суммарную поляризацию~\eqref{eq:polarization} можно пренебречь. В этом случае поперечное распределение поля моды определяется исключительно линейным откликом среды. Переходя в частотное представление, запишем волновое уравнение в спектральном представлении:

\begin{equation}\label{eq:helmholtz_vector}
	\nabla \times \nabla \times \widetilde{\mathbf{E}}(\mathbf{r},\omega) = \varepsilon(\omega) \frac{\omega^2}{c^2} \widetilde{\mathbf{E}}(\mathbf{r},\omega)
\end{equation}

Воспользуемся соотношением из векторного анализа

\begin{equation}\label{eq:rotrot_identity_delta}
	\nabla \times \nabla \times \widetilde{\mathbf{E}} = \nabla(\nabla \cdot \widetilde{\mathbf{E}}) - \Delta \widetilde{\mathbf{E}},
\end{equation}

В~\cref{eq:rotrot_identity_delta} сравним порядки величин первого и второго слагаемых.

\begin{subequations}\label{eq:e0_estimation}
	\begin{align}
		\Delta \widetilde{\mathbf{E}} &\sim k_0^2E_0, \\
		\nabla(\nabla \cdot \widetilde{\mathbf{E}}) &= \nabla(-\nabla\epsilon \cdot \widetilde{\mathbf{E}}/\epsilon) \sim \nabla(\dfrac{\Delta\epsilon E_0}{a \epsilon}) \sim \dfrac{\Delta \epsilon}{a^2 \epsilon}
	\end{align}
\end{subequations}

В~\cref{eq:e0_estimation} $E_0$~--- характерная величина амплитуды поля, $\delta$~--- толщина переходного слоя на границе <<сердцевина-оболочка>>, $\Delta \varepsilon$~--- перепад диэлектрической проницаемости между сердцевиной и оболочкой. Пренебрежение первым слагаемым допустимо при выполнении условия:

\begin{equation}\label{eq:weak_guidance_criterion}
	\frac{\Delta \varepsilon}{\varepsilon} \cdot \frac{1}{(k_0 \delta)^2} \ll 1
\end{equation}

Для типичного ОВ ($\delta \approx 1~\text{мкм}$, $\lambda \approx 1~\text{мкм}$, $\Delta n/n \sim 10^{-3}$) отношение $\frac{\Delta\varepsilon}{\varepsilon} \frac{1}{(k_0 \delta)^2}$ не превышает $10^{-4}$, что позволяет пренебречь слагаемым $\nabla(\nabla \cdot \widetilde{\mathbf{E}})$ и перейти к уравнению Гельмгольца:

\begin{equation}\label{eq:helmholtz_n}
	\nabla^2 \widetilde{\mathbf{E}}(\mathbf{r},\omega) + n^2(\omega) \frac{\omega^2}{c^2} \widetilde{\mathbf{E}}(\mathbf{r},\omega) = 0
\end{equation}

где поглощение в среде считается пренебрежимо малым, вследствие чего показатель преломления $n(\omega)$ полагается чисто вещественной величиной.
